% =====================================================================
%  Conferenze, Zagarolo (On The Road), giorno 2 — Talk 3
%  People Over Processes: il lavoro invisibile del producer
%  Fonte: trascrizione automatica "Track 17". Talk in inglese di un
%  producer e designer (Monument Valley 3, South of the Circle),
%  presentato da Claudia, con domande dal pubblico.
%  Etichette con prefisso "c3-3:". Nessun pacchetto aggiuntivo.
% =====================================================================

\chapter[People Over Processes]{People Over Processes: il lavoro invisibile del
producer}
\label{cap:c3-3-people-over-processes}

Il capitolo~\ref{cap:c2-2-producer} ha mostrato il lavoro del producer dal lato dei
processi: milestone, pianificazione, rapporti con il publisher. Questo talk, dal
titolo \textbf{People Over Processes}, ne mostra l'altra metà: il lavoro sulle
\textbf{persone}, spesso \textbf{invisibile}, che secondo il relatore è a volte il
più prezioso che un producer possa fare. Il capitolo è organizzato così:
\begin{enumerate}
  \item il relatore (sezione~\ref{sec:c3-3-relatore});
  \item perché è difficile parlare di produzione: contesto e invisibilità
        (sezione~\ref{sec:c3-3-difficile});
  \item il ciclo del problem solving collaborativo e le tre scale delle dinamiche
        umane (sezione~\ref{sec:c3-3-ciclo});
  \item le dinamiche uno a uno (sezione~\ref{sec:c3-3-uno-a-uno});
  \item le dinamiche di gruppo (sezione~\ref{sec:c3-3-gruppo});
  \item le dinamiche organizzative (sezione~\ref{sec:c3-3-organizzazione});
  \item il clima attuale e le conclusioni (sezione~\ref{sec:c3-3-conclusioni});
  \item le domande del pubblico (sezione~\ref{sec:c3-3-domande}).
\end{enumerate}

% ---------------------------------------------------------------------
\section{Il relatore}
\label{sec:c3-3-relatore}

Il relatore viene presentato da Claudia\footnote{Il nome del relatore non è
ricostruibile dalla trascrizione (``Red Cuddy''); l'introduzione lo presenta come
producer di uno studio trascritto ``Fast 2 Game'', con ogni probabilità
\textbf{ustwo games}, lo studio di \emph{Monument Valley}.}. Lavora nel settore
interattivo da quasi \textbf{vent'anni} (``spaventoso''), di cui circa
\textbf{quindici} specificamente nei videogiochi, ed è sia \textbf{designer} sia
\textbf{producer}.
\begin{description}
  \item[Come designer] ha progettato esperienze \textbf{VR per la Tate Modern},
        esperienze in \textbf{realtà aumentata} per promuovere serie TV e piccoli
        esperimenti: un flipper controllato con la faccia, un gioco in cui si
        ``ricostruisce'' la propria fidanzata in un robot, un altro in cui si è un
        adolescente\footnote{La descrizione è troncata nella trascrizione (``where
        you're a teenager with its protection'').}.
  \item[Come producer] ha realizzato \textbf{giochi educativi} per editori ed
        emittenti televisive; è stato \textbf{creative producer} dell'avventura
        narrativa \textbf{\emph{South of the Circle}} e, nel 2024, \textbf{lead
        producer} di \textbf{\emph{Monument Valley 3}}\footnote{Trascritto ``Money
        and Valley 3''.}, nello studio in cui lavora oggi.
\end{description}
Ha studiato alla \textbf{NFTS} (capitolo~\ref{cap:c2-4-nfts}). Avverte che
condivide volentieri il suo indirizzo email, ma ha una figlia di diciotto mesi ed è
così stanco che ``se una cosa non è nel mio campo visivo, smette di esistere'': non
risponderà.

% ---------------------------------------------------------------------
\section{Perché è difficile parlare di produzione}
\label{sec:c3-3-difficile}

\subsection{Il contesto: i cerchi concentrici}

Il relatore apre sempre i suoi talk con un'ammissione: il lavoro di produzione è
\textbf{molto difficile da raccontare}. La prima ragione è che dipende fortemente
dal \textbf{contesto}. Per spiegarlo usa una serie di \textbf{cerchi concentrici}:
\begin{enumerate}
  \item al centro c'è il \textbf{giocatore}, la persona più importante;
  \item intorno, ciò che sostiene il giocatore e ne soddisfa i bisogni: il
        \textbf{gioco};
  \item intorno al gioco, il \textbf{team} che lo fa;
  \item intorno al team, i \textbf{producer} e i \textbf{lead}: il lavoro
        manageriale, organizzativo e amministrativo;
  \item intorno ai producer, funzioni come le \textbf{risorse umane}.
\end{enumerate}
Esistono moltissimi tipi di giocatori, moltissimi giochi possibili per ciascuno e
moltissimi team possibili per ciascun gioco: le variabili si moltiplicano, e quando
si arriva al cerchio esterno il producer può dover essere in moltissimi modi diversi
per sostenere il proprio team.

\subsection{L'invisibilità}

La seconda ragione, che il relatore ha messo a fuoco di recente, è che il lavoro di
produzione è spesso \textbf{invisibile}:
\begin{itemize}
  \item \textbf{non è misurabile} in modo oggettivo;
  \item \textbf{non produce artefatti tangibili}: nessun file, nessuna animazione,
        nessun documento dimostra che il lavoro è stato fatto;
  \item è difficile dire \textbf{come} contribuisca al risultato: sappiamo che lo
        fa, ma non esattamente come.
\end{itemize}
Per questo il talk non propone regole universali ma un \textbf{quadro di
riferimento} per dare un nome a questo lavoro, che possa risuonare con le esperienze
di ciascuno, e che aiuti chi non ci ha mai fatto caso ad accorgersi che viene svolto
e a fargli spazio.

% ---------------------------------------------------------------------
\section{Il ciclo del problem solving e le tre scale}
\label{sec:c3-3-ciclo}

\subsection{Il viaggio e gli ostacoli}

Lo sviluppo di un gioco è collaborativo: persone con responsabilità, competenze e
priorità diverse lavorano insieme. Farle lavorare insieme \textbf{in modo efficace}
non è solo questione di processi, ma di capire, navigare e guidare le
\textbf{dinamiche umane}.

Immaginiamo di mettere insieme un team (o qualunque attività di gruppo): un po' di
soldi, un po' di hardware, uno spazio, ``il bagagliaio carico e pronti a mettersi
in strada''. Il viaggio fino al completamento è disseminato di
\textbf{ostacoli}, uno dopo l'altro.

\subsection{Il ciclo: identificare, discutere, decidere, agire}

Il modo in cui un gruppo risolve un problema si riduce, più o meno, sempre allo
stesso processo:
\begin{enumerate}
  \item \textbf{identificare} il problema;
  \item \textbf{discutere} insieme le soluzioni;
  \item \textbf{prendere una decisione};
  \item \textbf{agire}.
\end{enumerate}
L'azione genera a sua volta altri problemi, e il ciclo ricomincia. Il problema può
essere piccolissimo (``dove andiamo a pranzo?'') o enorme (``come pianifichiamo le
prossime otto settimane? Come raggiungiamo la prossima milestone?''), ma il ciclo è
sempre lo stesso.

\subsection{Tre scale}

Il relatore descrive le dinamiche umane su tre scale:
\begin{description}
  \item[Uno a uno.] Il rapporto tra due persone: un lead artist e il director, due
        artisti, due programmatori che si dividono il lavoro.
  \item[Gruppo.] La dinamica di un intero team, che richiede sensibilità diverse.
  \item[Organizzazione.] Il rapporto tra il team e un \textbf{organo di
        supervisione}: un leadership team, un consiglio di amministrazione, un
        publisher o, per uno studente, il docente o la commissione d'esame. Ne
        aveva parlato Manos il giorno prima (capitolo~\ref{cap:c2-6-miti-carriera}).
\end{description}
A ogni scala e a ogni fase del ciclo ci sono comportamenti che aiutano e
comportamenti che \textbf{bloccano}. Il compito del producer, il ``lavoratore
invisibile'', è riconoscere e rimuovere i comportamenti bloccanti il più in fretta
possibile.

% ---------------------------------------------------------------------
\section{Le dinamiche uno a uno}
\label{sec:c3-3-uno-a-uno}

\subsection{I comportamenti bloccanti}

Prendiamo un director e un lead artist. Ecco che cosa può andare storto in ogni
fase.
\begin{description}
  \item[Identificare.] Nessuno identifica il problema: il director dà feedback
        sulla direzione artistica, il lead artist cerca di accontentarlo, non
        funziona, ma nessuno dice ``c'è un problema''; si resta in un loop
        infinito. Oppure non tutti concordano che sia un problema: il lead artist
        dice ``siamo al settimo giro di feedback, così non si va avanti'', e l'altro
        non lo riconosce.
  \item[Discutere.] La discussione non è ammessa (``lo voglio verde, punto''); le
        proposte vengono respinte (``secondo me blu è meglio'' --- ``ma che dici?
        Fuori di qui''); oppure nessuno propone nulla: l'artista subisce
        passivamente (``dimmi cosa vuoi e lo faccio'').
  \item[Decidere.] La decisione non viene presa affatto: tante riunioni con molte
        opzioni discusse e poi ``ok, ottimo, ciao''. Oppure viene presa ma
        \textbf{non comunicata chiaramente} (``potrebbe essere verde, rosso o blu,
        mi piacciono tutti''), e ognuno esce con un'idea diversa: tempo sprecato.
        Oppure c'è \textbf{disaccordo dopo la decisione}: si è deciso verde, una
        settimana dopo è rosso.
\end{description}

\subsection{I segnali}

La scala uno a uno è la più vicina al modo in cui risolviamo i conflitti nella vita
privata, quindi i consigli sembrano ovvi, ma vale la pena dirli. Segnali da tenere
d'occhio:
\begin{itemize}
  \item \textbf{lamentele} in altre conversazioni (``non sopporto Dave, mi ha già
        fatto sette review e ancora non capisco che cosa vuole'');
  \item il \textbf{comportamento reciproco} nelle riunioni: il fastidio ``esce di
        traverso'' in altre occasioni;
  \item la \textbf{qualità del lavoro che cala}, perché non si arriva mai alla fase
        dell'azione.
\end{itemize}

\subsection{Le tecniche e i limiti}

\begin{itemize}
  \item Parlare con la persona che sembra avere il problema, prendendola da parte e
        \textbf{ascoltando apertamente e senza giudicare}.
  \item Se opportuno, parlare anche con l'altra persona, \textbf{senza riferire}
        ciò che ha detto la prima (``ho la sensazione che ci sia qualcosa che non
        va'').
  \item Accompagnare entrambi verso una soluzione e assicurarsi che sia una
        \textbf{soluzione chiara, riconosciuta da entrambi}; altrimenti ``si
        disperde nell'etere''.
\end{itemize}
Tre cautele:
\begin{enumerate}
  \item \textbf{mantenere un confine}: il producer non è un terapeuta; le
        conversazioni devono puntare a una soluzione;
  \item \textbf{evitare l'escalation}: c'è differenza tra ``mi dispiace, dev'essere
        difficile'' e ``che bastardo, non mi è mai piaciuto'';
  \item \textbf{riconoscere quando si è fuori dalla propria competenza}: è a questa
        scala che si incontrano più facilmente dinamiche problematiche (tutela delle
        persone, abusi, comportamenti inappropriati, problemi di salute,
        disabilità). In quei casi bisogna \textbf{passare subito il caso a chi è
        formato} per gestirlo.
\end{enumerate}

\subsection{Storia: la telecamera su binari}

Un caso vissuto dal relatore. Il team sviluppava un gioco 3D con una
\textbf{telecamera orbitale libera}, controllata dal giocatore. Il lead artist
propose al game director di tornare a una \textbf{telecamera su binari}, per
controllare meglio le composizioni. Il director: ``ne abbiamo parlato sei mesi fa,
abbiamo deciso, non si può''. Il lead artist: ``allora devo ridurre la qualità di
tutti gli asset, o mi serve più tempo''. Il director: ``ma gli obiettivi artistici
li hai fissati tu''. Muro contro muro.

Nel ciclo, la discussione non era ammessa e le proposte venivano respinte; il
problema era identificato, ma non tutti concordavano su quale fosse. Il producer
prese da parte il lead artist e scoprì che la vera preoccupazione era lo
\textbf{scope}: il team temeva la quantità di asset necessaria per un mondo
visibile da ogni angolazione. Così il lead artist era saltato direttamente a una
\textbf{soluzione} (la telecamera su binari, per costruire solo ciò che la camera
inquadra) invece di esporre il \textbf{problema} (``sono preoccupato per lo
scope''). Il director, davanti a quella che sembrava un \textbf{ultimatum}
(telecamera su binari o arte a bassa fedeltà), era andato nel panico. Una volta
capito questo, il producer chiese al lead artist di spiegare direttamente il
problema al director, aprendo la strada alla decisione e all'azione.

% ---------------------------------------------------------------------
\section{Le dinamiche di gruppo}
\label{sec:c3-3-gruppo}

\subsection{Che cosa cambia}

Due differenze importanti:
\begin{itemize}
  \item i comportamenti di gruppo \textbf{si rinforzano da soli}: se tre persone in
        una stanza fanno una cosa, la quarta tende ad adeguarsi;
  \item in un gruppo è facile pensare che \textbf{qualcun altro} sia responsabile.
\end{itemize}
Esempi: in riunione si percepisce che qualcosa non va, ma nessuno ne parla, quindi
``forse non è un problema''; il director si limita a dare ordini, e si capisce che
la discussione non è possibile; una decisione non è comunicata chiaramente, ma
nessuno si lamenta. Ed è molto più facile che \textbf{cinque} persone escano da una
conversazione con cinque idee diverse della decisione che due persone con due.

\subsection{Segnali e tecniche}

Segnali da osservare:
\begin{itemize}
  \item le \textbf{riunioni}, ma anche le interazioni su Slack o Discord;
  \item quando qualcuno solleva un problema: \textbf{viene risolto}? Spesso la
        conversazione va avanti, qualcuno apre una tangente e il problema iniziale
        non viene mai affrontato;
  \item i \textbf{problemi che ritornano}: il segno più sicuro che non vengono
        risolti.
\end{itemize}

Le conversazioni individuali, a questa scala, non bastano: quando si riporta a B
ciò che ha detto A, B aggiunge informazioni che bisognerebbe riportare ad A, e
intanto C\ldots{} Lo strumento principale è la \textbf{riunione di gruppo}.
\begin{itemize}
  \item \textbf{Essere disposti a dire la cosa scomoda.} Se ``l'atmosfera non va'',
        bisogna rompere il pensiero di gruppo e dire ciò che forse tutti pensano
        (``forse sono io, ma non è possibile che\ldots?''). Il relatore scherza
        sul fatto che sia ``una cosa molto inglese'' e che una sala di italiani
        forse non abbia questo problema.
  \item \textbf{Costruire una cultura}: se la cosa scomoda viene detta solo quando
        c'è il producer, il risultato non è ottimale. Bisogna rendere
        \textbf{sicuro} per gli altri dirla anche in sua assenza.
  \item \textbf{Dare l'esempio} quando la cosa scomoda riguarda proprio il
        producer: non mettersi sulla difensiva, ringraziare, rifletterci.
  \item Parlare in gruppo solo di \textbf{dinamiche generali}. Se il problema è
        \textbf{una persona}, non va affrontato davanti a tutti (``il tuo lavoro è
        carente da sei mesi''): non fa sentire nessuno al sicuro. Invece del
        \emph{calling out} (denunciare pubblicamente un comportamento), il relatore
        preferisce il \emph{calling in}: dopo la riunione prendere da parte la
        persona e parlarle.
\end{itemize}

\subsection{Storia: il team che non si criticava}

Dopo circa un anno di lavoro su un gioco, il team sembrava privo di
\textbf{entusiasmo} e di energia: le milestone e le review andavano bene, ma tutto
era ``fangoso'', lento. Nelle review nessuno diceva nulla di negativo, né sul gioco
né su altro. Il producer capì che tutti erano \textbf{molto disposti a sostenersi a
vicenda}: a ogni presentazione, ``bellissimo, bravo, grazie''. Ma \textbf{nessuno
criticava} il lavoro mostrato. Il ciclo si fermava alla prima fase: nessun problema
veniva identificato.

L'approccio fu doppio:
\begin{enumerate}
  \item una \textbf{raccolta di informazioni} individuale (``come ti senti riguardo
        al gioco? Va tutto bene?'');
  \item poi, in riunione, \textbf{prendere l'iniziativa}: ``non sono sicuro di
        quello che sta succedendo, ho la sensazione che\ldots''. Essere il primo a
        esporsi, e continuare a farlo finché gli altri non cominciano a farlo a loro
        volta.
\end{enumerate}

% ---------------------------------------------------------------------
\section{Le dinamiche organizzative}
\label{sec:c3-3-organizzazione}

\subsection{Che cosa cambia}

Gli \textbf{stakeholder} sono persone con autorità finanziaria e organizzativa sul
team: per uno studente il docente o la commissione, per uno studio un investitore,
un consiglio di amministrazione o un leadership team che gestisce più team. I
problemi del ciclo sono gli stessi, con due differenze:
\begin{itemize}
  \item quando succedono sono \textbf{dirompenti al massimo grado};
  \item richiedono \textbf{tantissimo tempo} per essere risolti, perché queste
        persone hanno pochissimo tempo, gestiscono molte situazioni simili e magari
        le si incontra una volta al mese.
\end{itemize}
Se un disaccordo tra art director e lead è una spia accesa sul cruscotto, un
problema con il consiglio di amministrazione è un \textbf{incidente a catena}:
tutto si ferma, tutto lo sforzo del team va nel risolverlo o nel dimostrare ciò che
va dimostrato, e quel tempo è tolto al gioco.

\subsection{Segnali}

Da osservare nella comunicazione con gli stakeholder, che di solito gestisce il
producer:
\begin{itemize}
  \item \textbf{ambiguità e vaghezza}: a quei livelli la cosa più sicura è spesso
        non dire nulla, o dire senza dire (``sì, credo che ci guarderemo'');
  \item il \textbf{panico}, o il bisogno di ``fare qualcosa'';
  \item l'\textbf{irrilevanza}: il tempo con queste persone è limitatissimo (magari
        un'ora al mese) ed è facile sprecarne dieci minuti su cose che non
        contano.
\end{itemize}

\subsection{Tecniche}

\begin{itemize}
  \item \textbf{Conoscere i cicli di reporting} dell'organizzazione: quanto spesso
        ci si incontra con il board, e quanto spesso il board deve riferire ai
        propri stakeholder. Se loro devono riferire a febbraio, bisogna dare loro
        qualcosa a gennaio. Se i budget vengono approvati ogni trimestre, perdere
        quella riunione significa aspettare tre mesi.
  \item \textbf{Costruire il consenso in anticipo}: con il tempo si costruiscono
        rapporti individuali con CEO, CTO, CMO; prima della riunione del lunedì si
        manda loro il materiale, ``così non è una sorpresa''. In quelle riunioni le
        decisioni raramente si prendono: \textbf{si confermano}.
  \item \textbf{Capire gli stakeholder}: che cosa li fa fare bella o brutta figura,
        che cosa vogliono, se preferiscono numeri o insistono per avere una demo.
        È conoscenza ``morbida'' che accelera l'approvazione.
  \item \textbf{Fare follow-up} il prima possibile, \textbf{per iscritto}, ``come
        se dovesse leggerlo un avvocato in tribunale'': è inutile ottenere una
        decisione se poi tutti dimenticano che cosa si intendeva.
\end{itemize}
Due cautele:
\begin{itemize}
  \item \textbf{niente sorprese in riunione}: in un ambiente ad alta tensione, dove
        nessuno vuole fare brutta figura e tutti vogliono giustificare il proprio
        ruolo, una sorpresa impedisce una risposta ponderata. Le reazioni saranno
        vaghe o drastiche (per esempio anticipare il lancio di tre mesi ``per
        mitigare il rischio'');
  \item \textbf{non accettare risposte ambigue}: se non si può avere una risposta
        nella riunione, si propone di approfondire a parte e si fa subito un
        follow-up scritto. L'ambiguità, a questo livello, ``è una campana a
        morto''.
\end{itemize}

\subsection{Storia: la vertical slice e il piano di produzione}

In uno studio in cui lavorava il relatore (con un altro producer responsabile del
processo), il leadership team doveva valutare una \textbf{vertical
slice}\footnote{Trascritta ``vertical site'', ``multiple slice'' e ``birth of the
size''.}, bellissima, su cui il team aveva lavorato moltissimo. Al gate andavano
presentati sia la vertical slice sia il \textbf{piano di produzione}, per ottenere
il budget dei mesi successivi e trasformarla in un gioco completo.

Il team, concentrato sulla vertical slice, preparò il piano di produzione
\textbf{solo la sera prima}, scoprendo che sarebbe durato \textbf{tre mesi in più}
del previsto. Fece comunque la riunione: un'ora e mezza di entusiasmo sulla
vertical slice, la riunione che sfora nell'ora di pranzo, e l'ultima slide con il
piano di tre mesi più lungo. ``Il gatto tra i piccioni'': nessuna decisione
possibile, domande affrettate (quanto costerà? si può risparmiare tempo?), tutto si
blocca.

La cosa giusta sarebbe stata fare per tempo una pianificazione seria. La seconda
cosa migliore: dire \textbf{all'inizio} della riunione ``vi mostriamo la vertical
slice, ma il piano di produzione non è pronto; dateci una settimana e ci rivediamo
al più presto per darvi queste informazioni critiche, fatte bene''.

% ---------------------------------------------------------------------
\section{Il clima attuale e le conclusioni}
\label{sec:c3-3-conclusioni}

\subsection{Dalla differenza tra bene e meglio a quella tra sopravvivere e
chiudere}

Dieci anni fa, dice il relatore, questo talk si sarebbe intitolato più o meno
``Il vostro team va bene? Ecco come farlo andare meglio''. Nel clima attuale questo
lavoro può fare la differenza \textbf{tra sopravvivere e chiudere}. Lo stress mette
sotto pressione ogni relazione: le persone hanno meno pazienza, i board sono più
inclini al panico, e i team sono stressati anche per altre ragioni.

Richiama le parole di apertura di Claudia: quando si è in viaggio, l'unica
risposta che si può ottenere è dalla persona alla propria sinistra o alla propria
destra. Questo lavoro invisibile, fatto di soft skill, è ciò che permette di
\textbf{parlare con chi ci sta accanto}.

\subsection{Sintesi}

\begin{itemize}
  \item Il lavoro di produzione è difficile da raccontare, e quello invisibile lo è
        ancora di più.
  \item Il lavoro invisibile consiste nel capire e navigare come le persone si
        comportano tra loro, come finiscono per comportarsi in gruppo e come deve
        comportarsi un'organizzazione per funzionare.
  \item Tutto ruota intorno al ciclo \textbf{identificare, discutere, decidere,
        agire}: quando funziona bene il ciclo scorre senza attriti.
\end{itemize}

\subsection{Come si sbaglia}

Il relatore ammette che un producer ha moltissimi modi per rovinare tutto. Per
esempio, se nessuno identifica i problemi, potrebbe creare un \textbf{foglio di
calcolo} in cui tutti scrivono i problemi, da rivedere ogni giorno in una lunga
riunione: oneroso, lento, impedisce alle persone di parlarsi, \textbf{sovra-processa}
ciò che dovrebbe essere una conversazione umana. Oppure stabilire che nessuna
decisione si prende senza di lui.

Il guaio del lavoro invisibile è che da fuori è difficilissimo distinguere quello
buono da quello cattivo: non ci sono misure oggettive. Nelle storie raccontate,
``sono stato io a far succedere qualcosa? Non lo so''. Ma questo lavoro fa la
differenza tra un producer che \textbf{aiuta} il team e uno che lo
\textbf{frena}. Il relatore cita un video in cui si diceva che i cattivi producer
``sono legione'': nei titoli di coda, soprattutto dei tripla A, si vedono team
gonfiati con moltissimi producer, scrum master e product manager che sembrano
contabili o amministratori (``c'è davvero così tanto da gestire?''). Ma ci sono
anche producer che \textbf{valgono oro}, e sono rarissimi.

% ---------------------------------------------------------------------
\section{Le domande del pubblico}
\label{sec:c3-3-domande}

\subsection{L'umiltà è la qualità chiave?}

Sì, ma un'umiltà particolare. Il relatore racconta che alla NFTS il responsabile del
corso di produzione cinematografica diceva:
\begin{quote}
l'ego del producer deve crescere e ridursi a comando.
\end{quote}
L'ego si \textbf{riduce} davanti al team, quando lo si guida: ``ragazzi, che cosa
stiamo facendo? Proviamo a farlo succedere'', tutti allo stesso livello. Si
\textbf{espande} quando il team è minacciato e quando bisogna ``dire la verità al
potere'': non in modo competitivo, ma per rappresentare e difendere i bisogni del
team quando il team non è nella stanza, a volte anche con forza.

Altre qualità:
\begin{itemize}
  \item \textbf{integrità}: ``la moglie di Cesare dev'essere al di sopra di ogni
        sospetto''. Il producer ha una posizione di potere unica (parla con il board
        e con il team, stabilisce il piano) e deve usarla con equità, mai a proprio
        vantaggio. Come diceva Manos il giorno prima, le persone non devono
        necessariamente \textbf{apprezzarti}, ma devono \textbf{fidarsi} di te e
        rispettarti;
  \item \textbf{concentrazione e chiarezza}: essere quello che dice ``aspettate,
        dieci minuti fa abbiamo lasciato in sospeso una cosa, possiamo tornarci?'';
  \item \textbf{coraggio}.
\end{itemize}

\subsection{Il producer che ricopre anche altri ruoli}

Nei piccoli team indie il producer è spesso anche director, modellatore o
designer. Come separare i ruoli? ``Con difficoltà'', e bisogna essere
\textbf{indulgenti con sé stessi}. Quando ci si chiede se un team abbia bisogno di
un producer, la risposta è: il team ha bisogno di \textbf{lavoro di produzione},
che non deve per forza farlo un producer. Indossare più cappelli fa parte del
lavoro indie. Consigli pratici:
\begin{itemize}
  \item dividere il tempo (``dopo pranzo sono il producer, prima no'');
  \item essere onesti con gli altri e con sé stessi nei momenti di conflitto: ``mi
        stai parlando come producer o come modellatore 3D?''; e chiederselo anche
        quando si risolve un problema;
  \item se aiuta, indossare \textbf{letteralmente} cappelli o magliette diverse per
        compartimentare;
  \item accettare che non sarà perfetto.
\end{itemize}

\subsection{Il producer deve proporre soluzioni?}

La priorità è la \textbf{risoluzione}: ogni conflitto che si trascina brucia tempo
e denaro. Proporre una soluzione va bene, purché lo si faccia con lo spirito giusto
(``la metto sul tavolo, vediamo se aiuta''). Il problema nasce quando ci si
\textbf{affeziona alla propria proposta}.

\subsection{I problemi irrisolvibili}

Quando si dice che un problema è irrisolvibile, spesso si intende che il gioco non
si può fare. Ma di solito significa che c'è un \textbf{punto fisso}: qualcosa che
non si può cambiare, e intorno a cui tutto il resto deve adattarsi. Non si può
creare tempo né ``far apparire soldi'', e ci sono persone e comportamenti che nessuna
conversazione cambierà. In quei casi l'attenzione va sulla
\textbf{produzione}, letteralmente: \textbf{aggirare} l'ostacolo e cambiare tutto il
resto per arrivare al traguardo. Il risultato forse sarà meno ambizioso del
previsto, ma il gioco esce. E allora bisogna dire agli stakeholder: questa è la
realtà, questa è la direzione; è accettabile per voi?

\subsection{Il lavoro da remoto}

Come identificare i problemi nei team che si vedono poco? Tornando in ufficio dopo il
Covid, il relatore ha notato una cosa: in presenza non si può fissare qualcuno che
non partecipa alla conversazione per vedere come reagisce, mentre in una
videochiamata si vedono tutti i volti. ``Ogni debolezza è una forza e ogni forza è
una debolezza.'' Da remoto si possono osservare le reazioni, fare riunioni molto
brevi, o iniziare una riunione con cinque persone e congedarle una alla volta man
mano che gli argomenti non le riguardano più. È difficile, ma bisogna aprire occhi
e mente per lasciar entrare quante più informazioni emotive e invisibili possibile.

\subsection{Il manager di più progetti}

Una persona che supervisiona più progetti, con pochissimo tempo, teme che i producer
pensino che abbia dei progetti preferiti. Che cosa può fare? Dal punto di vista di un
producer che riferisce a questa persona, il relatore vorrebbe:
\begin{itemize}
  \item \textbf{prove che abbia letto} i materiali inviati: domande sul contenuto
        della presentazione;
  \item che chieda ``come stai?'';
  \item un po' di tempo individuale, fuori dal team: magari \textbf{15 minuti a
        settimana} con ogni producer, ``un espresso insieme''.
\end{itemize}
``A volte è solo questione di gesti.''

\subsection{L'attrito è inevitabile?}

Anche con un ottimo producer, un po' di attrito nel team si accumulerà. Come
gestirlo? Il relatore non crede che si possa eliminare: i matrimoni più riusciti
sono quelli in cui ci si fa impazzire a vicenda di tanto in tanto. Creare una
cultura in cui l'attrito non esiste mai è ``una ricetta per il disastro''. Usa
un'immagine familiare:
\begin{quote}
va bene che mamma e papà litighino, purché i figli li vedano fare pace.
\end{quote}
Il conflitto in un team va bene purché tutti vedano che c'è una
\textbf{risoluzione}: se il litigio è pubblico, tutti devono vedere che è tutto a
posto; se è privato, entrambe le parti devono riconoscere che è risolto. Aspettare
che passi da solo non è un'opzione.

E se qualcuno odia segretamente un collega senza dirlo? Il relatore non ha una
tecnica segreta: ``mentirei''. In ogni impresa creativa bisogna
\textbf{aspirare alla perfezione sapendo di non raggiungerla}: se si punta al
grande si ottiene il buono, se si punta al buono si ottiene il discreto. Lo stesso
vale per la cultura del team: puntare a una salute perfetta, accettare di non
arrivarci, e continuare a spingere.

% ---------------------------------------------------------------------
\section*{Punti chiave del capitolo}
\addcontentsline{toc}{section}{Punti chiave del capitolo}

\begin{itemize}
  \item Il lavoro di produzione dipende dal \textbf{contesto} ed è spesso
        \textbf{invisibile}: non produce artefatti né misure oggettive.
  \item Ogni collaborazione ruota intorno al ciclo \textbf{identificare, discutere,
        decidere, agire}; il producer rimuove i comportamenti che lo bloccano.
  \item Uno a uno: ascoltare senza giudicare, ottenere una risoluzione chiara e
        condivisa, non fare il terapeuta, non alimentare l'escalation, passare i
        casi delicati a chi è formato. Spesso dietro una ``soluzione'' proposta c'è
        un problema non detto.
  \item Gruppo: dire la cosa scomoda e costruire una cultura in cui altri possano
        dirla; parlare in pubblico delle dinamiche, in privato delle persone
        (\emph{calling in}).
  \item Organizzazione: conoscere i cicli di reporting, costruire il consenso in
        anticipo, capire gli stakeholder, niente sorprese, niente ambiguità,
        follow-up scritto.
  \item Un buon producer ha un ego che cresce e si riduce a comando, integrità,
        chiarezza e coraggio; deve essere fidato più che amato.
  \item L'attrito è umano: conta che la risoluzione sia visibile.
\end{itemize}
